\section*{Bab \ref{ch:b2:waves-in-media} --- Gelombang Elektromagnetik dalam Plasma, Penghantar dan Dielektrik}

\begin{solution}{exo:b2:waves-in-media:1}
$f_p \approx 9\sqrt n$ Hz: \qty{2.8}{MHz}, \qty{9}{MHz}, \qty{9}{GHz}, \qty{2.6e15}{Hz}
($\lambda_p = \qty{115}{nm}$), \qty{2.8}{kHz}. Gelombang menengah pada \qty{1}{MHz}: hanya menembus
ruang antarbintang. FM: menembus ionosfer, bukan selubungnya. GPS:
segalanya kecuali selubungnya. Cahaya: segalanya kecuali tembaganya.
\end{solution}

\begin{solution}{exo:b2:waves-in-media:2}
Tembaga: \qty{9.2}{mm}, \qty{2.1}{mm}, \qty{65}{\micro m}, \qty{2.1}{\micro m}; aluminium pada
\qty{1}{MHz}: \qty{84}{\micro m}; air laut: \qty{2.3}{m} pada \qty{10}{kHz}, \qty{7}{mm} pada
\qty{1}{GHz}; tanah: \qty{5}{m} pada \qty{1}{MHz}. Kapal selam di \qty{10}{m} menuntut $\delta
\gtrsim \qty{10}{m}$: di bawah sekitar \qty{500}{Hz}.
\end{solution}

\begin{solution}{exo:b2:waves-in-media:3}
$n = 1.528$, $1.515$, $1.509$; $v = 1.963$, $1.980$, $\qty{1.987e8}{m/s}$; $n_g = n +
2B/\lambda^2 = 1.545$. Sudut biasnya $\arcsin(0.707/n)$: \qty{27.57}{\degree} dan
\qty{27.93}{\degree}: sebaran \qty{0.36}{\degree}.
\end{solution}

\begin{solution}{exo:b2:waves-in-media:4}
$f_p = \qty{4.0}{MHz}$: \qty{5}{MHz} lolos; \qty{2}{MHz} terpantul, menembus
$c/\sqrt{\omega_p^2 - \omega^2} = \qty{14}{m}$. Bagi $n = \qty{1.5e12}{m^{-3}}$: \qty{11}{MHz}.
Kerapatannya mengikuti fluks pengion Matahari --- jam, musim, dan
daur sebelas tahunnya --- sehingga frekuensi terpakai tertingginya bergerak bersamanya.
\end{solution}

\begin{solution}{exo:b2:waves-in-media:5}
(a) $1/\gamma\pi a^2 = \qty{5.3}{m\ohm/m}$. (b) $\delta = \qty{65}{\micro m}$; $R \approx 1/\gamma2\pi a\delta
= \qty{41}{m\ohm/m}$, $7.7$ kali yang searah. (c) $\delta = \qty{6.5}{\micro m}$, \qty{0.41}{\ohm/m},
$77$ kali. (d) $R \propto \sqrt f$ tumbuh secepat $L\omega$ begitu kulitnya
menguasai, dan tiap lilitan menambah hambatan; kawat Litz membundel banyak
untai terisolasi yang lebih tipis daripada $\delta$, sehingga seluruh tembaganya menghantar.
\end{solution}

\begin{solution}{exo:b2:waves-in-media:6}
$\delta_{\text{Al}} = 11.9$, $0.84$, \qty{0.084}{mm}: $d/\delta = 0.08$, $1.2$, $12$: $0.7$,
$10$, \qty{103}{dB}. Pada \qty{50}{Hz} penghantarnya bening; besi ber-$\mu_r$
tinggi mengalihkan fluksnya (perisaian magnetik) dan juga menyusutkan $\delta$ sebesar
$\sqrt{\mu_r}$.
\end{solution}

\begin{solution}{exo:b2:waves-in-media:7}
(a) $v_0 = eE_0/m\omega$; $\tfrac12nm\langle v^2\rangle = ne^2E_0^2/4m\omega^2 = \tfrac14\varepsilon_0
E_0^2\,\omega_p^2/\omega^2$. (b) $\tfrac14\varepsilon_0E_0^2$ dan $\tfrac14\varepsilon_0E_0^2(1 - \omega_p^2/
\omega^2)$. (c) Jumlahnya $\tfrac12\varepsilon_0E_0^2$ (kurungnya $1$); $\langle\Pi
\rangle = E_0^2k/2\mu_0\omega$, dan nisbahnya $c^2k/\omega = v_g$. (d) $\tfrac14\cdot\tfrac14
/\tfrac12 = 1/8$.
\end{solution}

\begin{solution}{exo:b2:waves-in-media:8}
(a) $1/v_g = (1/c)(1 - \omega_p^2/\omega^2)^{-1/2} \approx (1 + \omega_p^2/2\omega^2)/c$. (b) $\Delta t
= (L\omega_p^2/2c)(1/\omega_1^2 - 1/\omega_2^2)$ dengan $\omega_p^2 = ne^2/\varepsilon_0m$ dan $\omega =
2\pi f$: tetapannya $e^2/8\pi^2\varepsilon_0mc = \qty{1.34e-7}{}$ (SI). (c) $nL = \qty{9.3e23}{m^{-2}}$,
$(1/f_1^2 - 1/f_2^2) = \qty{5.7e-18}{s^2}$: $\Delta t = \qty{0.72}{s}$. (d) Kolomnya $nL$
(ukuran dispersinya), sehingga jaraknya bagi sebuah model $n$.
\end{solution}

\begin{solution}{exo:b2:waves-in-media:9}
(a) $\operatorname{Im}\underline\chi = \omega_p^2\Gamma\omega/[(\omega_0^2 - \omega^2)^2 + \Gamma^2\omega^2]$ dengan
$\omega_0^2 - \omega^2 \approx 2\omega_0(\omega_0 - \omega)$: $n'' = \tfrac12\operatorname{Im}\underline\chi$ adalah
Lorentzian yang dinyatakan itu. (b) $n''(\omega_0) = \omega_p^2/2\omega_0\Gamma$, $\alpha = 2n''\omega_0/c =
\omega_p^2/\Gamma c$. (c) $\omega_p^2 = \qty{3.2e20}{s^{-2}}$: $\alpha = \qty{1.8e4}{m^{-1}}$;
$\ln10/\alpha = \qty{0.13}{mm}$. (d) $n' - 1 \propto (\omega_0 - \omega)/[(\omega_0 - \omega)^2 +
\Gamma^2/4]$: positif di bawah, negatif di atas, menurun dalam $\pm\Gamma/2$
dari garisnya --- anomali di sana.
\end{solution}

\begin{solution}{exo:b2:waves-in-media:10}
(a) $\underline E = E_0\eu^{-x/\delta}\eu^{\iu(\omega t - x/\delta)}$, $\underline B = \underline k\underline E/\omega
= (1 - \iu)\underline E/\omega\delta$. (b) $\langle\Pi_x\rangle = \tfrac12\operatorname{Re}(\underline E
\underline B^*)/\mu_0 = E_0^2/2\mu_0\omega\delta$ di $x = 0$. (c) $\int_0^\infty\tfrac12\gamma E_0^2
\eu^{-2x/\delta}\dd x = \gamma E_0^2\delta/4 = E_0^2/2\mu_0\omega\delta$ dengan memakai $\gamma = 2/\mu_0\omega\delta^2$.
(d) $H_0 = \sqrt2E_0/\mu_0\omega\delta$, jadi $\tfrac12R_sH_0^2 = E_0^2/\gamma\mu_0^2\omega^2\delta^3 =
E_0^2/2\mu_0\omega\delta$. Tembaga pada \qty{1}{GHz}: $R_s = 1/\gamma\delta = \qty{8}{m\ohm}$;
$\tfrac12 \times 8 \times 10^{-3} \times 10^4 = \qty{40}{W/m^2}$.
\end{solution}

\begin{solution}{exo:b2:waves-in-media:11}
(a) $\omega_p = \qty{1.4e16}{rad/s}$, $\lambda_p = 2\pi c/\omega_p = \qty{140}{nm}$. (b) $\omega = \qty{3.8e15}{rad/s}
< \omega_p$: $1/\kappa = c/\sqrt{\omega_p^2 - \omega^2} = \qty{23}{nm}$, pemantulan total. (c)
Masih di bawah $\omega_p$: memantulkan; kebeningannya hanya di luar \qty{140}{nm} (pada
natrium, \qty{210}{nm}). (d) Tumbukan memberi bagian nyata pada $\underline\gamma$ dan
serapan beberapa persen; resonansi elektron terikat dalam cahaya tampak
(emas, tembaga) menyerap birunya lalu mewarnai logamnya.
\end{solution}

\begin{solution}{exo:b2:waves-in-media:12}
(a) $\omega\tau = 0.12$: $\underline\varepsilon_r = 84 - 9.7\iu$; pada \qty{20}{GHz}, $\omega\tau = 1$:
$45 - 40\iu$. (b) $\underline n \approx 9.2 - 0.53\iu$; $c/2n''\omega = \qty{1.8}{cm}$: ovennya
memasak dua sentimeter terluarnya, hantaran mengerjakan sisanya. (c) $\tfrac12 \times
1.54 \times 10^{10} \times 8.85 \times 10^{-12} \times 9.7 \times 4 \times 10^6 = \qty{2.6}{MW/m^3}$. (d) Pada
\qty{20}{GHz} panjang serapannya akan semilimeter: hanya kulitnya
yang akan matang; \qty{2.45}{GHz} menembus, dan merupakan pita bebas.
\end{solution}

\begin{solution}{pb:b2:waves-in-media:1}
\textbf{1.} $f_p = \qty{9}{MHz}$: GPS lolos, gelombang pendeknya terpantul.

\textbf{2.} $k = (\omega/c)\sqrt{1 - \omega_p^2/\omega^2}$: $v_\varphi = \omega/k \approx c(1 + \omega_p^2/
2\omega^2)$, $v_g = \dd\omega/\dd k = c^2k/\omega \approx c(1 - \omega_p^2/2\omega^2)$.

\textbf{3.} $\Delta t = \int(1/v_g - 1/c)\dd z = \int\omega_p^2\dd z/2c\omega^2 = e^2N_T/2\varepsilon_0
mc\omega^2 = (e^2/8\pi^2\varepsilon_0mc)N_T/f^2 = 1.34 \times 10^{-7} \times 2 \times 10^{17}/2.48 \times
10^{18} = \qty{11}{ns}$: \qty{3.2}{m}.

\textbf{4.} $N_T = (\Delta t_1 - \Delta t_2)/K(1/f_1^2 - 1/f_2^2)$; selisihnya
\qty{7}{ns}: $1\%$ menuntut \qty{0.07}{ns}.

\textbf{5.} $v_\varphi$ melampaui $c$ sebanyak $v_g$ kurang darinya: fase
pembawanya tiba lebih awal sementara modulasinya tiba terlambat. Penerima
berfase pembawa menerapkan koreksinya dengan tanda yang berlawanan.

\textbf{6.} \qty{3.2}{m} pada siang hari, \qty{0.3}{m} pada malam hari.

\textbf{7.} $f_p = \qty{0.9}{MHz}$ dengan tumbukan: gelombang menengah dan pendek
terserap (sebuah pemadaman radio); GPS pada \qty{1.5}{GHz} merasakan serapan
$\propto 1/\omega^2$, yang dapat diabaikan.

\textbf{8.} $f_1\Delta t = 1.575 \times 10^9 \times 1.1 \times 10^{-8} = 17$ daur.

\textbf{9.} $\delta = \sqrt{2/\mu_0\gamma\omega} = \qty{10}{\micro m}$: dinding semilimeter adalah
seratus \omterm{def:b2:dispersion-wave-packets:complex}{kedalaman kulit}.

\textbf{10.} $R_s = 1/\gamma\delta = \qty{0.1}{\ohm}$; $\tfrac12R_sH_0^2 = \qty{45}{W/m^2}$: \qty{45}{W},
$6\%$.

\textbf{11.} $\eu^{-\pi \times 0.5} = 0.21$: \qty{-14}{dB} dalam amplitudo, \qty{-27}{dB} dalam
daya (dan lubang mungilnya memang sedikit sekali berkopel sejak awal).

\textbf{12.} Cahaya, $\lambda = \qty{0.5}{\micro m}$, dua ribu kali lebih kecil
daripada lubangnya lalu lolos; gelombang mikronya, $\lambda = \qty{12}{cm}$,
seratus kali lebih besar.

\textbf{13.} $\underline n = 2 - 0.005\iu$; $c/2n''\omega = \qty{1.9}{m}$: kacanya
nyaris tak menghangat.

\textbf{14.} $j \approx H_0/\delta = \qty{3e7}{A/m^2}$, $j^2/\gamma = \qty{9e8}{W/m^3}$ dalam
lapisan sepuluh mikrometer: ujungnya memanas dalam hitungan milisekon, melontarkan elektron,
mengionkan udaranya --- bunga api.

\textbf{15.} $\delta = \qty{18}{\micro m}$; penampangnya $2\pi a\delta = \qty{5.6e-8}{m^2}$ melawan
$\pi a^2 = \qty{7.9e-7}{m^2}$: $R_{\text{ac}} = \qty{0.30}{\ohm/m}$, empat belas kali yang searah.

\textbf{16.} $RI^2 = \qty{120}{W/m}$; permukaannya $\pi d = \qty{3.1e-3}{m^2}$ tiap meter:
$\Delta T \approx \qty{1900}{K}$ --- ia akan mencair; kumparan semacam itu berupa tabung berpendingin air.

\textbf{17.} $\delta = \qty{0.14}{mm}$: arus imbasnya, dan kalornya, tinggal
dalam sepersepuluh milimeter; sebuah pulsa pendek yang disusul pencelupan mengeraskan
permukaan sebuah roda gigi sementara intinya tetap liat.

\textbf{18.} $R = 1/\gamma2\pi a\delta = \qty{0.10}{\ohm/m}$, tiga kali lebih kecil;
bagian dalamnya tak mengemban apa pun --- dari situlah tabungnya, dengan pendinginnya di tempat
tembaganya akan terbuang percuma.

\textbf{19.} $\delta = a$ pada $f = 1/\pi\mu_0\gamma a^2 = \qty{17}{Hz}$; di bawahnya, arusnya
seragam dan rumus searahnya berlaku.

\textbf{20.} $n^2 \approx 1 + \omega_p^2/\omega_0^2 + (\omega_p^2/\omega_0^2)(\lambda_0/\lambda)^2$: yaitu bentuk Cauchy
dengan $n(\infty)^2 - 1 = \omega_p^2/\omega_0^2 = 1.25$ bagi $n(\infty) = 1.5$.

\textbf{21.} $n = 1.531$ dan $1.510$, $v = 1.96$ dan \qty{1.99e8}{m/s}; $n_g(550) =
1.550$; $n_g(400) - n_g(700) = 0.063$: \qty{210}{ns} sepanjang satu kilometer.

\textbf{22.} $\sin((A + D)/2) = 1.5165 \times 0.5$: $D = \qty{38.6}{\degree}$; $\dd D = 2\sin
(A/2)\dd n/\cos((A + D)/2) = 1.53\,\dd n$: $\Delta D = \qty{1.85}{\degree}$, \qty{6.5}{cm} pada \qty{2}{m}.

\textbf{23.} Pada \qty{200}{nm} frekuensinya mendekati resonansinya dan
$n''$ tumbuh: buram. Resonansi lemah pada \qty{700}{nm} menyerap merahnya:
kaca hijau.

\textbf{24.} Biru lebih dekat ke resonansi ultraungunya: elektron
terikatnya menanggapi lebih kuat, $n$ lebih besar, dan simpangannya pun.

\textbf{25.} Lapisan F pada \qty{1.5}{GHz}: $\underline\gamma$ khayal, perambatan
dengan tunda grup yang kecil. Baja pada \qty{2.45}{GHz}: $\underline\gamma$ nyata, \omterm{thm:b2:waves-in-media:skin}{efek
kulit}, pemantulan. Tembaga pada \qty{13.56}{MHz}: hal yang sama, dengan kulit
\qty{18}{\micro m}. Kaca dalam cahaya tampak: $\iu\omega\varepsilon_0\chi$, indeks nyata, sebuah
\omterm{def:b2:dispersion-wave-packets:relation}{medium dispersif} yang bening.
\end{solution}
